\documentclass{standalone}
\usepackage{standalone}

\begin{document}
\chapter{Limitations and Future Work}
\label{ch:limitations}

\section{Current Limitations}
\begin{itemize}
  \item \textbf{Technician payouts are manual.} The wallet keeps exact track of earnings and withdrawals, but the money itself is sent by an admin through bKash Send Money, who then records the TrxID. There is no automatic payout through a bKash, Nagad or bank API.
  \item \textbf{Card payments run in test mode.} Stripe does not offer local accounts to businesses in Bangladesh. A real launch would need a local gateway such as SSLCOMMERZ or the bKash Payment Gateway, plugged in behind the same payment service interface.
  \item \textbf{Pickup capacity is checked but not locked.} If two customers book the last place in a slot in the same second, both bookings could go through. At the current scale this is acceptable, since admins can see every pickup time, but a busy service would need an atomic slot counter.
  \item \textbf{Prices are provisional.} The 45 services and their price ranges are a starting point for demonstration, not prices agreed with a real repair business. Admins can change them at any time.
  \item \textbf{A technician can have only one active job.} This prevents double booking, but it does not fit technicians who handle several small jobs in a day.
  \item \textbf{Some internal values are left over from the old codebase.} Stored values such as \texttt{rider} and \texttt{parcel\_delivered} come from the parcel-delivery codebase. Users never see them, but a migration would make the data model cleaner.
  \item \textbf{The browser tests do not cover signed-in screens.} They only cover the public pages, because signing in needs real Firebase accounts. The signed-in screens are covered instead by server tests of every endpoint behind them, and I checked them by hand (Chapter~\ref{ch:walkthrough}).
  \item \textbf{There are no real customers yet.} The platform is deployed and tested, but it has not yet been used by real customers and technicians with real money, and it has not had a load test or an external security audit.
\end{itemize}

\section{Future Enhancements}
\begin{enumerate}
  \item \textbf{A local payment gateway and automatic payouts}: SSLCOMMERZ or bKash for both customer payments and technician payouts, with reconciliation.
  \item \textbf{Live technician location}: showing the technician's approach on a map on the day of the pickup.
  \item \textbf{A mobile app or Progressive Web App}: a technician app that works well offline and supports push notifications.
  \item \textbf{SMS notifications}: many customers in Bangladesh rely on SMS more than on email.
  \item \textbf{Parts inventory}: recording the parts used in each repair, with their cost and supplier.
  \item \textbf{Technician scheduling}: daily capacity, working hours and assignment that takes travel routes into account.
  \item \textbf{Cash payment}: letting customers pay in cash at hand-over, with the platform's share settled from the technician's wallet.
  \item \textbf{A proper domain name}: a \texttt{.com} or \texttt{.com.bd} address, since some mobile networks could not resolve the free \texttt{.vercel.app} address.
  \item \textbf{A Bangla interface}: translating the whole interface, not just accepting Bangla input for the AI.
  \item \textbf{Cleaning up the schema}: removing the leftover values through a planned migration.
\end{enumerate}

\end{document}
